\chapter{Requirements Specification} \label{cap:requirements}

This chapter defines the requirements for the eight GCC cross-compiler targets developed in this work.
Each target is a modified GCC backend that accepts bare-metal freestanding C programs and produces
machine code containing only the instructions supported by the corresponding educational processor.

\section{Scope}

All targets compile \textbf{bare-metal freestanding C}: programs with no standard library and no
operating system, producing an ELF binary as output.

\section{Correctness Requirements}

Two independent correctness requirements apply to every target.

\textbf{ISA compliance}: the compiler must never emit an instruction whose opcode is not in the
target's allowed set. This requirement is verifiable by inspecting the assembly output.

\textbf{Behavioral equivalence}: the compiled program must exhibit the same observable behavior as
the source --- the same return values, the same memory effects, and the same control flow --- as if
compiled for a processor that natively supports the full instruction set.

\section{rvsc0 --- Basic Single-Cycle}

The rvsc0 target corresponds to the eight-instruction single-cycle processor described in Chapter 4.4
of \cite{patterson2020}, \textit{A Simple Implementation Scheme}:

\begin{quote}
In this section, we look at what might be thought of as a simple implementation of our RISC-V subset.
[\ldots] This simple implementation covers load word (lw), store word (sw), branch if equal (beq),
and the arithmetic-logical instructions add, sub, and, and or.
\end{quote}

\begin{table}[htb]
\centering
\caption{rvsc0 native instruction set}
\label{tab:rvsc0-isa}
\begin{tabular}{ll}
\toprule
\textbf{Mnemonic} & \textbf{Description} \\
\midrule
\texttt{lw}   & Load word \\
\texttt{sw}   & Store word \\
\texttt{beq}  & Branch if equal \\
\texttt{add}  & Add \\
\texttt{addi} & Add immediate \\
\texttt{sub}  & Subtract \\
\texttt{and}  & Bitwise AND \\
\texttt{or}   & Bitwise OR \\
\bottomrule
\end{tabular}
\end{table}

\subsection{Limitations}

The rvsc0 processor has no mechanism to write the program counter --- it supports neither \texttt{jalr}
nor any jump-and-link instruction. As a consequence, subroutine calls are architecturally impossible,
and the set of C programs that can be compiled for rvsc0 is restricted accordingly. For all other C
constructs, the compiler must produce correct behavior using only the native instruction set.

\section{rvsc1 --- Extended Single-Cycle}

The rvsc1 target extends rvsc0 with \texttt{lui} and \texttt{jalr}. These two instructions are the
minimum addition that enables the full C calling convention: \texttt{lui} materializes 32-bit absolute
addresses and \texttt{jalr} performs indirect jumps and encodes function returns via
\texttt{jalr x0, 0(ra)}.

\begin{table}[htb]
\centering
\caption{rvsc1 native instruction set}
\label{tab:rvsc1-isa}
\begin{tabular}{ll}
\toprule
\textbf{Mnemonic} & \textbf{Description} \\
\midrule
\texttt{lw}   & Load word \\
\texttt{sw}   & Store word \\
\texttt{beq}  & Branch if equal \\
\texttt{add}  & Add \\
\texttt{addi} & Add immediate \\
\texttt{sub}  & Subtract \\
\texttt{and}  & Bitwise AND \\
\texttt{or}   & Bitwise OR \\
\texttt{lui}  & Load upper immediate \\
\texttt{jalr} & Jump and link register \\
\bottomrule
\end{tabular}
\end{table}

With these ten instructions, the target must support the full C calling convention and general-purpose
C programs.

\section{rvsc2 --- Single-Cycle Without Fence and Control}

The rvsc2 target implements the full RV32I base ISA with three groups of instructions excluded. Memory
ordering instructions are unnecessary on a single-core single-cycle processor: omitting them is
behavior-preserving. CSR access and system instructions require operating system support that is absent
in the bare-metal environment.

\begin{table}[htb]
\centering
\caption{rvsc2 native instruction set (RV32I base, excluding fence, CSR, and system groups)}
\label{tab:rvsc2-isa}
\begin{tabular}{ll}
\toprule
\textbf{Group} & \textbf{Instructions} \\
\midrule
Integer arithmetic & \texttt{add, sub, addi, and, andi, or, ori, xor, xori, sll, slli, srl, srli, sra, srai} \\
Comparison         & \texttt{slt, sltu, slti, sltiu} \\
Loads              & \texttt{lw, lh, lhu, lb, lbu} \\
Stores             & \texttt{sw, sh, sb} \\
Branches           & \texttt{beq, bne, blt, bge, bltu, bgeu} \\
Jumps              & \texttt{jal, jalr} \\
Upper immediate    & \texttt{lui, auipc} \\
\bottomrule
\end{tabular}
\end{table}

\begin{table}[htb]
\centering
\caption{rvsc2 instruction treatment}
\label{tab:rvsc2-treatment}
\begin{tabular}{ll}
\toprule
\textbf{Treatment} & \textbf{Instructions} \\
\midrule
Behaviorally equivalent & \texttt{fence, fence.i, sfence.vma} \\
Rejected & \texttt{ecall, ebreak, sret, wfi}, CSR instructions (\texttt{csrrw, csrrs, csrrc, csrrwi, csrrsi, csrrci}) \\
\bottomrule
\end{tabular}
\end{table}

\section{rvsc3 --- Single-Cycle rv32i}

The rvsc3 target implements the complete RV32I base integer ISA \cite{riscv-spec}, including the
memory ordering, CSR access, and system instruction groups omitted from rvsc2. No instruction
synthesis is required: every RV32I instruction is natively supported by the hardware.

\begin{table}[htb]
\centering
\caption{rvsc3 native instruction set (complete RV32I)}
\label{tab:rvsc3-isa}
\begin{tabular}{ll}
\toprule
\textbf{Group} & \textbf{Instructions} \\
\midrule
Integer arithmetic & \texttt{add, sub, addi, and, andi, or, ori, xor, xori, sll, slli, srl, srli, sra, srai} \\
Comparison         & \texttt{slt, sltu, slti, sltiu} \\
Loads              & \texttt{lw, lh, lhu, lb, lbu} \\
Stores             & \texttt{sw, sh, sb} \\
Branches           & \texttt{beq, bne, blt, bge, bltu, bgeu} \\
Jumps              & \texttt{jal, jalr} \\
Upper immediate    & \texttt{lui, auipc} \\
Memory ordering    & \texttt{fence, fence.i} \\
System             & \texttt{ecall, ebreak, wfi, sret} \\
CSR                & \texttt{csrrw, csrrs, csrrc, csrrwi, csrrsi, csrrci} \\
\bottomrule
\end{tabular}
\end{table}

\section{rvsc4 --- Single-Cycle rv64i}

The rvsc4 target is a 64-bit processor implementing the full rv64i instruction set. In addition to all
rvsc3 instructions, it natively supports:

\begin{itemize}
  \item 64-bit memory operations: \texttt{ld}, \texttt{sd}, \texttt{lwu}
  \item Word operations with sign-extended 64-bit results: \texttt{addw}, \texttt{subw},
        \texttt{addiw}, \texttt{sllw}, \texttt{srlw}, \texttt{sraw}, \texttt{slliw},
        \texttt{srliw}, \texttt{sraiw}
\end{itemize}

The W-suffix variants operate on 32 bits and sign-extend the result to 64 bits, following the same
convention as \texttt{addw}, \texttt{subw}, and the other W instructions from rvsc4.

\section{rvsc5 --- Single-Cycle rv64im}

The rvsc5 target extends rvsc4 with the M extension, adding native integer multiply and divide:

\begin{itemize}
  \item Multiply: \texttt{mul, mulw, mulh, mulhu, mulhsu}
  \item Signed divide: \texttt{div, divw}
  \item Unsigned divide: \texttt{divu, divuw}
  \item Signed remainder: \texttt{rem, remw}
  \item Unsigned remainder: \texttt{remu, remuw}
\end{itemize}

\section{rvsc6 --- Single-Cycle Floating-Point}

The rvsc6 target extends rvsc5 with the complete F (single-precision) and D (double-precision)
floating-point extensions as defined in \cite{riscv-spec}, supporting all arithmetic, memory,
conversion, and comparison instructions for both precisions.

\begin{itemize}
  \item Loads/stores: \texttt{flw, fsw} (single); \texttt{fld, fsd} (double)
  \item Arithmetic: \texttt{fadd.s, fadd.d, fsub.s, fsub.d, fmul.s, fmul.d, fdiv.s, fdiv.d,
        fsqrt.s, fsqrt.d}
  \item Fused multiply-add: \texttt{fmadd.s, fmadd.d, fmsub.s, fmsub.d, fnmadd.s, fnmadd.d,
        fnmsub.s, fnmsub.d}
  \item Sign, minimum, maximum: \texttt{fsgnj.s, fsgnj.d, fsgnjn.s, fsgnjn.d, fsgnjx.s, fsgnjx.d,
        fmin.s, fmin.d, fmax.s, fmax.d}
  \item Comparison and classification: \texttt{feq.s, feq.d, flt.s, flt.d, fle.s, fle.d,
        fclass.s, fclass.d}
  \item Move between integer and FP registers: \texttt{fmv.w.x, fmv.x.w} (single);
        \texttt{fmv.d.x, fmv.x.d} (double)
  \item Conversions (rv64): \texttt{fcvt.s.w, fcvt.s.wu, fcvt.s.l, fcvt.s.lu} (int to single);
        \texttt{fcvt.w.s, fcvt.wu.s, fcvt.l.s, fcvt.lu.s} (single to int);
        \texttt{fcvt.d.w, fcvt.d.wu, fcvt.d.l, fcvt.d.lu} (int to double);
        \texttt{fcvt.w.d, fcvt.wu.d, fcvt.l.d, fcvt.lu.d} (double to int);
        \texttt{fcvt.s.d, fcvt.d.s} (between single and double)
\end{itemize}

\section{rvsc7 --- Single-Cycle Atomic}

The rvsc7 target extends rvsc6 with the A extension (atomic memory operations), implementing the full
rv64imafd instruction set \cite{riscv-spec}. The \texttt{.w} variants operate on 32 bits
(sign-extended to 64); the \texttt{.d} variants operate on 64 bits.

\begin{itemize}
  \item Load-reserved and store-conditional: \texttt{lr.w, lr.d}; \texttt{sc.w, sc.d}
  \item AMO (atomic memory operations):
        \texttt{amoadd.w, amoadd.d} (atomic add);
        \texttt{amoand.w, amoand.d} (atomic AND);
        \texttt{amoor.w, amoor.d} (atomic OR);
        \texttt{amoxor.w, amoxor.d} (atomic XOR);
        \texttt{amoswap.w, amoswap.d} (atomic swap);
        \texttt{amomax.w, amomax.d} (atomic signed maximum);
        \texttt{amomaxu.w, amomaxu.d} (atomic unsigned maximum);
        \texttt{amomin.w, amomin.d} (atomic signed minimum);
        \texttt{amominu.w, amominu.d} (atomic unsigned minimum)
\end{itemize}
